\documentclass[11pt]{article}
\usepackage[margin=1in]{geometry}
\usepackage{amsmath}
\usepackage{booktabs}
\usepackage{url}
\title{Joint estimator--controller design for the noisy-sensor encounters}
\date{Prepared October 5, 2026}
\begin{document}
\maketitle

\section{Question}
The user asked for the estimator and the controller to be treated as one
design, and for the encounters with the noisy synthetic sensors to work.
Working here means no collision, a solved command in every frame, and
recovery to the given path. The user left the design open. This report
records the coupled design, the measured results and the development steps
that led to it.

\section{Starting point (aec20bc)}
With exact target and ego states, all 14 encounters recovered. With the NRMM
estimator and noisy sensors (seed 20261003), 5 of 14 recovered and none
collided; the other 9 ended without a solved command
(\texttt{pcbfNoNumericalResult}). The estimator's ego bound was invalid in 9
of 13 runs, because the maneuvers left its declared domain
(\path{ESTIMATOR_VALIDITY_20261004.tex}). Where the bound was published, it
failed to contain the truth in most 15-m/s frames: the kinematic premise
$v_y=l_Er$ of the velocity measurement does not hold while the rear tire
slips (\path{ESTIMATOR_DOMAIN_20261005.tex}). The ego heading error reached
0.07--0.31~rad (95th percentile) and the lateral velocity error 1--3.9~m/s at
15~m/s.

\section{Design}
The changes come in pairs: what the controller provides to the estimator,
and what the estimator provides to the controller.

\paragraph{E1. Shared model and input.} The estimator receives the input the
controller held and the controller's vehicle and Fiala tire model. The
lateral force balance $m a_y=F_{x,f}\sin\delta+F_{y,f}\cos\delta+F_{y,r}$
replaces $v_y=l_Er$. Each axle force is monotone in $v_y$, so the
accelerometer determines $v_y$, and a certified interval follows from the
sensor bounds, a 5\% box on stiffness and friction, and the sideslip cone.
The published value is the nominal solution, not the interval midpoint. The
domain (speed, yaw rate, acceleration, yaw acceleration, sideslip cone) is
computed from the controller's limits and tire friction.
The velocity bound now includes the drift of the held measurement between
samples (\path{estimator/MODEL_AIDED_ESTIMATION.md}).

\paragraph{E2. Contract-consistent target parameters.} The target contract
keeps $A$ and $\beta$ constant. A Gauss--Newton fit of the closed-form NRMM
trajectory to the last 2~s of radar detections replaces the tracker's $A$ and
$\beta$ in the published estimate. The detections are converted to the
inertial frame with headings obtained by integrating the gyro back from the
current yaw estimate. $A$ and the curvature are kept only where they exceed
three standard errors. Until the window spans 0.5~s, $A=\beta=0$ is published
instead of the tracker's initial transient. The tracker's velocity, position
and certified sets are kept, and the published radii are widened by the
distance between the two values.

\paragraph{E3. Acquisition and initialization.} The 0.5-s radar window before
$t=0$ also initializes the tracker's acceleration: the window is fitted with a
constant acceleration, and only components above three standard errors are
kept. The window is converted with the GNSS course fitted linearly over it;
one noisy course sample rotates a 20-m detection by about 0.12~m. A target
that enters the radar range during a run is acquired after
0.25~s of detections, whose least-squares line initializes its velocity. A
single detection with the 15-m/s line-of-sight prior had started a re-acquired
8-m/s target 7~m/s off.

\paragraph{C1. Information-consistent holds.} Every hold of the plan starts
from the estimator's current enclosure and propagates it through that hold
only. The target set restarts from its predicted state at the hold start.
The open-loop image over the variable horizon (up to 512 holds) had made the
hard road and state rows infeasible.

\paragraph{C2. Stable-handling envelope.} At both ends of every hold the plan
keeps the rear tire below Fiala saturation under the hold's braking ratio,
$\|[(v_y-l_rr)/k;\ \bar v_xb]\|\le v_x$, on the estimate. It also keeps the
sideslip inside the estimator's cone, robust to the enclosure. This removes
the spin observed at 15~m/s and keeps the force-balance measurement
informative. The adhesion cone is scaled to unit entries.

\paragraph{C3. CLF on the estimate.} The CLF row $V(\hat x^+)\le\rho
V(\hat x)$ uses the estimate. The worst case over the enclosure cannot be met
near the path, and minimizing its slack made the braking ratio chatter.

\paragraph{C4. Slack structure with an enclosure.} The untightened collision
rows are hard only after the slack prefix, as without an enclosure. The
tightened rows carry slack at every node.

\paragraph{C5. Road recovery in the terminal set.} The endpoint must be able
to stop its lateral motion toward either road edge before reaching it:
$\max(w_e,0)^2\le2a\,d_e$ with $a=0.5\mu g$, as one rotated cone per edge.
The horizon rollout stops only where this holds with half the deceleration.

\paragraph{C6. Startup seed.} The potential-field seed counts the ego's own
heading in its clearance envelope. It chooses the passing side by the road
room that remains after clearing the target, rather than by the sign of a
noisy sideslip estimate.

\section{Results}
\paragraph{Encounters.} Each encounter ran for up to 2000 holds (100~s), with
recovery confirmed after the last given-path collision time and a 1-s dwell
inside the tolerances. On the circular road the given path meets the target
again near 80~s, so curved runs confirm recovery after that second encounter.
\begin{center}
\begin{tabular}{lrrrrr}
\toprule
 & runs & recovered & no solution & collision & min.\ gap (m)\\
\midrule
exact states, aec20bc & 14 & 14 & 0 & 0 & 0.124\\
exact states, this design & 14 & 14 & 0 & 0 & 0.153\\
noisy, seed 20261003, aec20bc & 14 & 5 & 9 & 0 & 0.161\\
noisy, six seeds, this design & 84 & 84 & 0 & 0 & 0.416\\
\bottomrule
\end{tabular}
\end{center}
The noisy seeds are 20261003 and 1--5. No run left the road. The smallest road
margin was 1.00~m (noisy) and 3.85~m (exact). The median of the noisy minimum
gaps was 1.64~m, against 0.20~m with exact states. The controller plans
against a forecast it knows to be uncertain, and its tightened rows keep more
clearance. Recovery on the straight road was confirmed after a median of
10.9~s (noisy) and 10.7~s (exact), at most 13.5~s
(\path{JOINT_DESIGN_20261005/campaign-metrics.csv}).

\paragraph{Estimation.} The 95th percentile of the ego estimation error per
run, over the 84 noisy runs, was as follows:
\begin{center}
\begin{tabular}{lrr}
\toprule
 & aec20bc (seed 20261003) & this design (six seeds)\\
\midrule
heading, 8~m/s & 0.017--0.068 rad & 0.0015--0.0061 rad\\
heading, 15~m/s & 0.069--0.31 rad & 0.0023--0.0065 rad\\
lateral velocity, 8~m/s & 0.12--0.47 m/s & 0.019--0.059 m/s\\
lateral velocity, 15~m/s & 0.98--3.85 m/s & 0.009--0.106 m/s\\
position & & 0.007--0.022 m\\
\bottomrule
\end{tabular}
\end{center}
The published target parameters were compared with the truth on the
seed-20261003 runs (root mean square per run, median over 14 runs). From
0.5~s to 3~s, the error of $A$ fell from 0.27~m/s$^2$ (tracker, stage F) to
0.021~m/s$^2$, and that of $\beta$ from 0.0066 to 0.0006~rad. In the first
0.5~s the largest $\beta$ error fell from 0.044 to 0.0077~rad. The $A$ error
there stays about 0.3~m/s$^2$: model selection holds $A$ at zero until it is
significant.

\paragraph{Truth audit.} \path{scripts/replayEstimatorValidity.m} replayed the
estimator on every frame of the 14 seed-20261003 runs and reproduced the
published estimates to $5\times10^{-15}$. The ego bound was valid in every
frame. Neither the published ego bounds nor the target bounds failed to
contain the truth in any frame where they were available (5,896 ego and 2,247
target frames).
The true state left the shared domain in one frame: in the 15-m/s
head-on the rear-adhesion ratio was 1.008 at 1.25~s, because the adhesion
row acts on the estimate and on the affine prediction
(\path{JOINT_DESIGN_20261005/audit.csv}). For comparison, at aec20bc the ego
bound was lost in 9 of 13 runs, and at 15~m/s the available bounds failed to
contain the truth in 2--24 frames per run.

\paragraph{Computation.} The four exact encounters 8-m/s turning crossing,
8-m/s head-on, 15-m/s crossing and 15-m/s head-on were run one at a time on the
same machine, about 800 holds each way. The controller took a median of
9.2~ms per hold (aec20bc: 6.7~ms) and a 95th percentile of 41~ms (30~ms). It
exceeded the 50-ms hold in 3.4\% of holds (1.4\%). At equal horizon length the
cost rose by about 1.4--1.5 times, because of the adhesion and road-recovery
cones (\path{JOINT_DESIGN_20261005/timing-exact-sequential.csv}). With the
estimator, the observer took a median of 11--15~ms per hold depending on
machine load (aec20bc: 8.6~ms, run alone). In the final campaign, run as eight
parallel processes, the controller exceeded 50~ms in 7.1\% of the noisy
holds. Deadline misses are recorded but are not a failure criterion of these
experiments.

Two changes are neutral for the trajectories, which were verified
bit-identical on two noisy runs. A cache of the cruise trim never hit,
because \texttt{strcmp} does not match MATLAB string objects inside a cell
array, so every hold re-solved the trim. The per-hold target tube is now
evaluated in one vectorized call. Splitting the adhesion cones to hold ends
only saved about 5\% and was not adopted.

\paragraph{Development record.} The full noisy campaign (seed 20261003) after
each stage:
\begin{center}
\begin{tabular}{llr}
\toprule
stage & added & recovered\\
\midrule
aec20bc & --- & 5/14\\
A & E1 (force balance, shared domain) & 2/14\\
B & C1 (per-hold enclosure) & 4/14\\
F & C3, C2 (adhesion robust), E3 (acceleration fit) & 10/14\\
I & C4, C5, C6, E3 significance test, velocity-bound drift & 11/14\\
J & E2 (window fit, per-detection headings) & 14/14\\
L & cone scaling & 14/14\\
final & adhesion on estimate, model selection, gyro headings, & 14/14\\
 & $A=\beta=0$ until 0.5~s, 0.25-s acquisition & \\
\bottomrule
\end{tabular}
\end{center}
Stage A lost three runs: the bounds became valid, the open-loop tube entered
the hard rows, and the target tube drove large swerves toward the road
edge. In stage B, four runs did not recover within 100~s because the robust
CLF chattered. At stage L, seeds 1--3 gave 14, 13 and 14; the failure was the
robust adhesion reserve. Two later stages were stopped early. With the adhesion
row on the estimate, two 15-m/s head-on runs of seed 4 had no solution in
their first frame: a noisy curvature of the straight target caused it. With
model selection added, the 8-m/s curved head-on of seed 5 had no solution
when the target was re-acquired for its second encounter. These led to the
final changes (\path{JOINT_DESIGN_20261005/development-metrics.csv}).

\section{Scope and limits}
\begin{itemize}
\item The plant in these simulations is the controller's model with nominal
  parameters, so the force-balance certificate's structural premise holds by
  construction; the 5\% parameter box is a design margin, not a measured
  mismatch.
\item Only the first hold of each plan uses a certified enclosure. Later holds
  assume that the estimator's enclosure at their start equals the present one.
  The target's forecast error beyond one hold is not in the constraints; the
  plan relies on the window fit and on replanning.
\item The window fit is an estimate, not an enclosure; the certified target
  parameter set remains wide (course about $\pm0.4$~rad, the full contract
  range for $A$ and curvature).
\item The rear-adhesion row acts on the estimate and is not part of the safety
  certificate. The road-recovery deceleration fraction (0.5) is a declared
  setting.
\item Recursive feasibility is not claimed; every frame either returns a
  solved command or reports that there is none.
\item During this work MATLAB sessions began to end with a heap-corruption
  message at exit, after all outputs had been written. A session running only
  \texttt{disp(1+1)} showed it too, so the message does not come from the
  project code. It appeared in the run during which the user quota of the
  RAM-backed \path{/tmp} filled, and it disappeared once \texttt{TMPDIR}
  pointed to disk. Orphaned MATLAB service processes from those sessions
  then made new sessions fail to start (error 5001). The final campaign ran
  with \texttt{TMPDIR} on disk, after those processes were stopped.
\end{itemize}

\section{Addendum: E2 removed (October 5, 2026)}
At the user's direction the window fit (E2) was removed and the tracker's own
$A$ and $\beta$ are published again; everything else above is unchanged.
With the tracker's output the six-seed noisy campaign recovered 82 of 84
encounters, with no collision or road departure. In the other two, the
controller had no solution where a long horizon (84 and 104 holds) met a
forecast moved by the tracker's error in $A$. The smallest road margin
fell to 0.23~m. The truth audit again found no bound that failed to contain
the truth (\path{ESTIMATOR_DIRECT_TARGET_20261005.tex}).

\section{Reproduction}
Each encounter was run in its own MATLAB process with the arguments that
\texttt{runPotentialFieldCampaign} passes
(\path{JOINT_DESIGN_20261005/campaignOne.m.txt}):
\begin{verbatim}
c = estimatorControllerIntegrationConfig(); c.randomSeed = seed;  % noisy
runNonlinearPredictiveSafetyValidation(Scenarios=name, Frames=2000, ...
  ControllerConfiguration=struct('referenceSpeed',speed, ...
    'controller',struct('horizonSteps',8+8*(speed==15))), ...
  RequireCollisionThreat=true, RecoveryDwellSeconds=1, ...
  EstimatorConfiguration=c)          % struct() for exact states
\end{verbatim}
For one seed this is \texttt{runPotentialFieldCampaign(root, out, 2000,
EstimatorConfiguration=c)}. The truth audit is
\texttt{replayEstimatorValidity(out)} on the seed-20261003 outputs. The
estimator's ego kernel must be rebuilt after this change
(\path{scripts/buildNrmmObserverKernel.m}); generated binaries under
\path{solver/} are not part of the change. Outputs are under
\path{~/.cache/collisionAvoidance/joint-design-20261005}.
\end{document}
